\begin{table}[htbp]
\centering
\begin{minipage}{14.9058cm}
\caption{\label{tab:syr2_mr_comparison_k2} MR and MCL violation rates and counts by SYR2 contaminant reporting status, two-step upstream watershed linkage, 1985--2005}
\end{minipage}
\small
{\setlength{\tabcolsep}{4pt}%
\begin{tabular}{>{\raggedright\arraybackslash}p{6.5cm}>{\raggedleft\arraybackslash}p{1.75cm}>{\raggedleft\arraybackslash}p{1.75cm}>{\raggedleft\arraybackslash}p{1.75cm}>{\raggedleft\arraybackslash}p{1.75cm}}
\toprule
 & \multicolumn{2}{c}{Utility appears in SYR2 (\%)} & & \\
\cmidrule(lr){2-3}
 & \multicolumn{1}{>{\centering\arraybackslash}p{1.75cm}}{Yes} & \multicolumn{1}{>{\centering\arraybackslash}p{1.75cm}}{No} & \multicolumn{1}{>{\centering\arraybackslash}p{1.75cm}}{Diff.} & \multicolumn{1}{>{\centering\arraybackslash}p{1.75cm}}{$p$-value} \\
 & \multicolumn{1}{>{\centering\arraybackslash}p{1.75cm}}{(1)} & \multicolumn{1}{>{\centering\arraybackslash}p{1.75cm}}{(2)} & \multicolumn{1}{>{\centering\arraybackslash}p{1.75cm}}{(3)} & \multicolumn{1}{>{\centering\arraybackslash}p{1.75cm}}{(4)} \\
\midrule
Ever MR violation & 51.27\phantom{$^{***}$} & 44.08\phantom{$^{***}$} & 7.19$^{*\phantom{*}\phantom{*}}$ & 0.09\phantom{$^{***}$} \\
Ever MCL violation & 5.51\phantom{$^{***}$} & 3.55\phantom{$^{***}$} & 1.96$^{\phantom{*}\phantom{*}\phantom{*}}$ & 0.28\phantom{$^{***}$} \\
Mean annual violations (count, 0--6) & 0.13\phantom{$^{***}$} & 0.20\phantom{$^{***}$} & \makebox[0pt][r]{-}0.07$^{**\phantom{*}}$ & 0.01\phantom{$^{***}$} \\
\bottomrule
\end{tabular}
}
\begin{minipage}{14.9058cm}
\vspace{4pt}
\footnotesize
\raggedright
\textit{Notes:} Sample restricted to utilities with a coal mine within two flow steps upstream of their intake and no coal mine colocated with their intake. \textit{Has SYR2 reading}: at least one Six-Year Review contaminant reading for arsenic, nitrate, selenium, or barium. \textit{No SYR2 reading}: utility is in a state where at least one utility has a SYR2 reading but the utility itself does not. MR violation: inorganic chemical monitoring and reporting violation. MCL violation: inorganic chemical maximum contaminant violation. Mean annual violations: sum of MR and MCL violation indicators across arsenic, nitrate, and inorganic chemicals categories (maximum of 6 per utility-year). $N$ utilities: group 1 = 236, group 2 = 338. The first two rows are percentages of utilities and their differences are in percentage points. Stars on differences are from two-sample $t$-tests. *** p$<$0.01, ** p$<$0.05, * p$<$0.1.
\end{minipage}
\end{table}
